% Options for packages loaded elsewhere
% Options for packages loaded elsewhere
\PassOptionsToPackage{unicode}{hyperref}
\PassOptionsToPackage{hyphens}{url}
\PassOptionsToPackage{dvipsnames,svgnames,x11names}{xcolor}
%
\documentclass[
  letterpaper,
  DIV=11,
  numbers=noendperiod]{scrartcl}
\usepackage{xcolor}
\usepackage{amsmath,amssymb}
\setcounter{secnumdepth}{-\maxdimen} % remove section numbering
\usepackage{iftex}
\ifPDFTeX
  \usepackage[T1]{fontenc}
  \usepackage[utf8]{inputenc}
  \usepackage{textcomp} % provide euro and other symbols
\else % if luatex or xetex
  \usepackage{unicode-math} % this also loads fontspec
  \defaultfontfeatures{Scale=MatchLowercase}
  \defaultfontfeatures[\rmfamily]{Ligatures=TeX,Scale=1}
\fi
\usepackage{lmodern}
\ifPDFTeX\else
  % xetex/luatex font selection
\fi
% Use upquote if available, for straight quotes in verbatim environments
\IfFileExists{upquote.sty}{\usepackage{upquote}}{}
\IfFileExists{microtype.sty}{% use microtype if available
  \usepackage[]{microtype}
  \UseMicrotypeSet[protrusion]{basicmath} % disable protrusion for tt fonts
}{}
\makeatletter
\@ifundefined{KOMAClassName}{% if non-KOMA class
  \IfFileExists{parskip.sty}{%
    \usepackage{parskip}
  }{% else
    \setlength{\parindent}{0pt}
    \setlength{\parskip}{6pt plus 2pt minus 1pt}}
}{% if KOMA class
  \KOMAoptions{parskip=half}}
\makeatother
% Make \paragraph and \subparagraph free-standing
\makeatletter
\ifx\paragraph\undefined\else
  \let\oldparagraph\paragraph
  \renewcommand{\paragraph}{
    \@ifstar
      \xxxParagraphStar
      \xxxParagraphNoStar
  }
  \newcommand{\xxxParagraphStar}[1]{\oldparagraph*{#1}\mbox{}}
  \newcommand{\xxxParagraphNoStar}[1]{\oldparagraph{#1}\mbox{}}
\fi
\ifx\subparagraph\undefined\else
  \let\oldsubparagraph\subparagraph
  \renewcommand{\subparagraph}{
    \@ifstar
      \xxxSubParagraphStar
      \xxxSubParagraphNoStar
  }
  \newcommand{\xxxSubParagraphStar}[1]{\oldsubparagraph*{#1}\mbox{}}
  \newcommand{\xxxSubParagraphNoStar}[1]{\oldsubparagraph{#1}\mbox{}}
\fi
\makeatother


\usepackage{longtable,booktabs,array}
\usepackage{calc} % for calculating minipage widths
% Correct order of tables after \paragraph or \subparagraph
\usepackage{etoolbox}
\makeatletter
\patchcmd\longtable{\par}{\if@noskipsec\mbox{}\fi\par}{}{}
\makeatother
% Allow footnotes in longtable head/foot
\IfFileExists{footnotehyper.sty}{\usepackage{footnotehyper}}{\usepackage{footnote}}
\makesavenoteenv{longtable}
\usepackage{graphicx}
\makeatletter
\newsavebox\pandoc@box
\newcommand*\pandocbounded[1]{% scales image to fit in text height/width
  \sbox\pandoc@box{#1}%
  \Gscale@div\@tempa{\textheight}{\dimexpr\ht\pandoc@box+\dp\pandoc@box\relax}%
  \Gscale@div\@tempb{\linewidth}{\wd\pandoc@box}%
  \ifdim\@tempb\p@<\@tempa\p@\let\@tempa\@tempb\fi% select the smaller of both
  \ifdim\@tempa\p@<\p@\scalebox{\@tempa}{\usebox\pandoc@box}%
  \else\usebox{\pandoc@box}%
  \fi%
}
% Set default figure placement to htbp
\def\fps@figure{htbp}
\makeatother





\setlength{\emergencystretch}{3em} % prevent overfull lines

\providecommand{\tightlist}{%
  \setlength{\itemsep}{0pt}\setlength{\parskip}{0pt}}



 


\KOMAoption{captions}{tableheading}
\makeatletter
\@ifpackageloaded{caption}{}{\usepackage{caption}}
\AtBeginDocument{%
\ifdefined\contentsname
  \renewcommand*\contentsname{Table of contents}
\else
  \newcommand\contentsname{Table of contents}
\fi
\ifdefined\listfigurename
  \renewcommand*\listfigurename{List of Figures}
\else
  \newcommand\listfigurename{List of Figures}
\fi
\ifdefined\listtablename
  \renewcommand*\listtablename{List of Tables}
\else
  \newcommand\listtablename{List of Tables}
\fi
\ifdefined\figurename
  \renewcommand*\figurename{Figure}
\else
  \newcommand\figurename{Figure}
\fi
\ifdefined\tablename
  \renewcommand*\tablename{Table}
\else
  \newcommand\tablename{Table}
\fi
}
\@ifpackageloaded{float}{}{\usepackage{float}}
\floatstyle{ruled}
\@ifundefined{c@chapter}{\newfloat{codelisting}{h}{lop}}{\newfloat{codelisting}{h}{lop}[chapter]}
\floatname{codelisting}{Listing}
\newcommand*\listoflistings{\listof{codelisting}{List of Listings}}
\makeatother
\makeatletter
\makeatother
\makeatletter
\@ifpackageloaded{caption}{}{\usepackage{caption}}
\@ifpackageloaded{subcaption}{}{\usepackage{subcaption}}
\makeatother
\usepackage{bookmark}
\IfFileExists{xurl.sty}{\usepackage{xurl}}{} % add URL line breaks if available
\urlstyle{same}
\hypersetup{
  pdftitle={Maior Faturamento},
  colorlinks=true,
  linkcolor={blue},
  filecolor={Maroon},
  citecolor={Blue},
  urlcolor={Blue},
  pdfcreator={LaTeX via pandoc}}


\title{Maior Faturamento}
\author{}
\date{}
\begin{document}
\maketitle


\subsection{Maior Faturamento}\label{maior-faturamento}

Boa tarde, chefe, tudo bem?

Estive trabalhando em uma possibilidade para aumentar o faturamento da
empresa. Com isso, notei quatro pontos que valem a pena serem
discutidos:

\textbf{1º} A categoria dos eletrônicos, embora não venda em mesma
quantidade que os acessórios, gera um faturamento muito superior

\begin{longtable}[]{@{}lll@{}}
\toprule\noalign{}
Categoria & Faturamento & Vendas \\
\midrule\noalign{}
\endhead
\bottomrule\noalign{}
\endlastfoot
Eletrônicos & 79.300 & 41 \\
Acessórios & 13.000 & 60 \\
\end{longtable}

\textbf{2º} Os produtos que mais geram faturamento são o celular e o
tablet, com \textbf{36.800} e \textbf{27.100}, respectivamente. No
entanto, o celular não é oferecido em ambas as cidades, somente em São
Paulo. Até por isso, acredito que os números entre ambas as cidades
sejam tão discrepantes, com a metrópole paulista (\textbf{61.000})
faturando, quase \textbf{2 vezes mais} que o Rio de Janeiro
(\textbf{31.300}).

\textbf{3º} Em relação ao tablet, temos as seguintes informações:

\begin{verbatim}
# A tibble: 2 x 5
# Groups:   produto [1]
  produto cidade preco_medio total_vendas faturamento
  <chr>   <chr>        <dbl>        <dbl>       <dbl>
1 Tablet  RJ            1450           11       15950
2 Tablet  SP            1600            7       11200
\end{verbatim}

Ou seja, adotando um preço reduzido, como no RJ, obtivemos um número
maior de vendas, resultando em um \textbf{faturamento 42\% maior} que em
SP. \textbf{11.200} para \textbf{15.950} (aumento de 42\%).

\textbf{4º} Já em relação ao celular, produto que gera maior
faturamento, tomando por base uma previsão feita por mim, é provável que
suas vendas \textbf{não tenham o mesmo comportamento que as vendas do
tablet}. Logo, caso o preço de celular seja reduzido, ele, provavelmente
não trará faturamento maior. A análise que usei para constatar essas
informações tem valor estatisticamente seguro.

\begin{verbatim}
       1 
9.497788 
\end{verbatim}

Esse é o número aproximado de vendas que cada um dos 2 modelos de
celular teriam caso seus preços fossem \textbf{1.900}, ou seja,
arredondando o valor para 9 e calculando o novo faturamento total dos
celulares, teríamos o valor de \textbf{34.200}, menor que o atual
(\textbf{36.800}).

\subsection{O que fazer?}\label{o-que-fazer}

Por tudo isso, minhas sugestões para aumentar o faturamento total de
nossa empresa se resumem também a 4 ideias principais:

\begin{itemize}
\item
  Investir na produção de mais tipos de eletrônicos, pois retornam
  grande faturamento.
\item
  Leve reajuste no preço do tablet na cidade de São Paulo, assim o suas
  vendas podem seguir o comportamento das vendas no Rio de Janeiro e,
  com isso, aumentar o faturamento.
\item
  Iniciar a venda de celulares no Rio de Janeiro, porque é o produto que
  gera maior faturamento, isso pode inclusive equilibrar o faturamento
  gerado entre as duas cidades.
\item
  Não reduzir o preço dos celulares, ao menos por hora, uma vez que isso
  potencialmente não traria benefícios
\end{itemize}

Desde já, agradeço

Caio Augusto




\end{document}
